In this paper, we describe Data Flyway, an in-flight processing system that enables HPC applications to easily create and run in-flight transformation and analysis pipelines. Our implementation aims to optimize the computation of data transformations and analysis operations and data movement across multiple tiers of the HPC memory and storage hierarchy. Our evaluation demonstrates that Data Flyway removes transformation cost from the application's critical path in the VPIC-IO workload by overlapping transformations with computation, achieving $\sim$3--7$\times$ lower observed I/O time than HDF5 with the asynchronous VOL connector. Our dynamic transformation scheduler reduces total transformation time by 7--22\% compared to static scheduling, saving hundreds of seconds at scale, and for iterative ML workloads Data Flyway achieves approximately 55\% higher I/O throughput than HDF5 at all scales by serving pre-transformed data from a server-side cache shared across nodes. For analysis, performing a single-stage vector magnitude pipeline in the data path is $\sim$2$\times$ faster than a post-hoc PDC workflow and $\sim$15$\times$ faster than a post-hoc HDF5 workflow, and a multi-stage curl-vorticity pipeline with transient intermediates requires $\sim$30\% less time than its post-hoc equivalent at every scale. By performing transformations and analysis closer to the storage system and overlapping computation with data movement, Data Flyway reduces unnecessary data transfer and improves overall system efficiency. In future work, we plan to extend the analysis graph to express spatial data dependencies, allowing stencil-based analyses such as curl to exchange halo regions between PDC servers so that boundary values are computed exactly.

\newpage